% Imporditud: tools/import_eksperimendid.py
% Allikas: Eesti füüsikaolümpiaadi eksperimendi ülesannete
%   kogu (Rael Kalda, 2023), ülesanne Ü5, lahendus L5.
\setAuthor{EFO žürii}
\setRound{piirkonnavoor}
\setYear{1998}
\setNumber{G E1}
\setDifficulty{1}
\setTopic{staatika}
\setVahendid{joonlaud, pliiats, tuntud massiga koormis või münt}
\setPunktid{8}
\setAllikas{Eksperimendi kogu Ü5, lahendus lk 35}
\setLahendusseis{toores}

\prob{Joonlaua mass}
Määrake joonlaua mass. Hinnake mõõtmistulemuse täpsust.

\hint

\solu
Teooria: Joonlaud, mille ühe otsa peale on asetatud münt, tasakaalustatakse pliiatsil. Joonlaud on tasakaalus, kui seda päri- ja vastupäeva pööravad jõumomendid on suuruselt võrdsed. Jõumoment arvutatakse valemist M = mgL, kus m on keha mass, g vaba langemise kiirendus, L jõuõlg (1 p). Tähistame $L_1$ joonlaua selle osa pikkuse, kus asub münt ja $L_2$ joonlaua teise osa pikkuse. Kummalegi poole toetuspunkti jääb osa joonlauast. Homogeense massijaotuse korral, võime joonlaua massi joontiheduse avaldada seosest $\rho$ = m/L, kus L = $L_1$ + $L_2$ (1 p). Sel juhul on joonlaua mass ühel pool toetuspunkti võrdne m = $\rho$$L_1$ ja teisel pool toetuspunkti m = $\rho$$L_2$. Jõumomentide seos saab kuju m1gL1 + $\rho$L1gL1/2 = $\rho$L2g $\cdot$ $L_2$/2, kus tähega $m_1$ on tähistatud mündi mass (1 p).

Tuletatud seosest saame avaldada joonlaua joontiheduse

$\rho$ = (2m1L1) / [($L_2$ $-$ $L_1$) $\cdot$ ($L_2$ + $L_1$)]

ja sellest omakorda joonlaua massi m = $\rho$L. Suhtelise vea valem: E(x) = $\Delta$x/x. Summa (vahe) arvutamise viga: E ($x_1$ $\pm$ $x_2$) = ($\Delta$$x_1$ + $\Delta$$x_2$) / ($x_1$ $\pm$ $x_2$) Korrutise (jagatise) arvutamise viga: E ($x_1$/$x_2$) = E ($x_1$ $\cdot$ $x_2$) = E ($x_1$)+E ($x_2$) (1 p).

Mõõtmised: Paigutame mündi joonlaua ühe otsa peale. Tasakaalustame mündiga joonlaua pliiatsil. Mõõdame joonlaua kaugused $L_1$ ja $L_2$ (1 p).

Arvutused: Arvutame eeltoodud valemite abil joonlaua joontiheduse ja massi (1 p).

Mõõtmisvea hindamine: Arvutame massi suhtelise vea E(m) = $\Delta\rho$/$\rho$ + $\Delta$L/L (1 p). $\Delta$L/L võime jätta arvestamata, kuna see on suhteliselt väike. Seega E(m) = $\Delta\rho$/$\rho$. Joontiheduse vea arvutamisel võime jätta arvestamata ka liikme ($L_2$ + $L_1$) vea, kuna ka see on väike võrreldes liikme ($L_2$ $-$ $L_1$) veaga, seega

$\Delta\rho$/$\rho$ = $\Delta$$L_1$/$L_1$ + ($\Delta$$L_1$ + $\Delta$$L_2$) / ($L_2$ $-$ $L_1$) (1 p).
\probend
